\documentclass[letterpaper, 10 pt, conference]{ieeeconf}

\IEEEoverridecommandlockouts
\overrideIEEEmargins

\title{\LARGE \bf
Robot Localization for Triton Using Particle Filters
}

\author{Shivraj Bhatti}

\begin{document}
\maketitle
\thispagestyle{empty}
\pagestyle{empty}

\begin{abstract}
This report presents a Monte Carlo localization system for a simulated Triton
robot in a known indoor map. The system combines an odometry-based motion model
with a likelihood-field LiDAR sensor model and uses particle filtering to
estimate robot pose over time. Placeholder text will be replaced with measured
results after the Gazebo experiments are complete.
\end{abstract}

\section{Introduction}
Robot localization estimates the pose of a robot relative to a known map. In
this project, Triton localizes in a Gazebo house environment using a particle
filter. The filter represents belief as weighted pose samples and repeatedly
performs prediction, correction, and resampling.

\section{Approach}
\subsection{Map and Likelihood Field}
The map was generated using ROS GMapping and saved as an occupancy-grid map.
The likelihood field was computed as a distance transform over occupied cells,
so each LiDAR beam endpoint can be scored by distance to the nearest obstacle.

\subsection{Motion Model}
The motion model uses consecutive odometry messages to compute the
rotation-translation-rotation decomposition. Gaussian noise is applied to each
motion component before propagating particles.

\subsection{Sensor Model}
The sensor model evaluates a subsampled LaserScan. For each valid beam, the
beam endpoint is transformed into the map frame and scored with a zero-centered
Gaussian over nearest-obstacle distance from the likelihood field.

\subsection{Particle Filter}
Particles are initialized uniformly in free map cells with random yaw. Each
iteration applies odometry prediction, LiDAR likelihood weighting, weight
normalization, pose estimation, and low-variance resampling.

\section{Experiments}
TODO: Insert map and likelihood-field figures.

TODO: Insert qualitative screenshots showing particle convergence over time.

TODO: Insert quantitative position and yaw error plots from
\texttt{artifacts/pf\_metrics.png}.

\section{Results and Discussion}
TODO: Summarize convergence behavior, final error, failure cases, and sensitivity
to map distortion, beam subsampling, and particle count.

\section{Conclusion}
TODO: State whether the robot localized successfully and describe future
improvements such as adaptive resampling, random particle injection, and stronger
handling of kidnapped-robot cases.

\end{document}
